\documentclass[11pt]{article}
\usepackage[margin=1in]{geometry}
\usepackage{amsmath,amssymb,amsthm}
\usepackage{graphicx}
\usepackage[hidelinks]{hyperref}

\newtheorem{theorem}{Theorem}
\newtheorem{lemma}{Lemma}
\newtheorem{remark}{Remark}

\newcommand{\F}{\mathcal F}
\newcommand{\Lip}{\operatorname{Lip}}
\newcommand{\R}{\mathbb R}
\newcommand{\N}{\mathbb N}

\title{A positive answer for Kaufmann's cusp question}
\author{FA/Banach automatic math run}
\date{2026-06-30}

\begin{document}
\maketitle

\section{Source question}

The source paper is P. L. Kaufmann, \emph{Products of Lipschitz-free spaces and applications}, arXiv:1403.6605, Studia Math. 226 (2015), 213--227.  After Proposition 5.1, the paper asks:

\begin{quote}
Consider the set
\[
 Cusp=\{(x,0):x\geq 0\}\cup \{(x,x^2):x\geq 0\}
\]
endowed with the Euclidean metric. Is \(\F(Cusp)\) isomorphic to a subspace of \(L^1\)?
\end{quote}

\begin{center}
\includegraphics[width=0.95\linewidth]{figures/open_problem_crop.png}
\end{center}

\section{Proof intuition}

The obstruction in the source question is the tangency at the origin: the two branches are not orthogonal there, so Kaufmann's direct union formula does not apply.  The tangency can nevertheless be organized scale by scale.

Place short rungs between the two branches at the points \(1/n\) near the origin and at the integers far away.  The Euclidean cusp is bi-Lipschitz equivalent to the two rails with the induced metric from this weighted ladder graph.  The rung spacing is chosen so that near the origin the horizontal spacing is comparable with the cusp thickness \(x^2\), while at infinity integer spacing is already fine because the cusp thickness is large.

The full ladder is a metric graph.  Its free space is an \(L^1\)-edge space modulo the graph's cycle space.  Choosing the small upper edges and the large vertical rungs as chords gives fundamental cycles whose \(L^1\)-lengths are uniformly controlled by their chord lengths.  Therefore the cycle space is complemented in the edge \(L^1\) space.  The quotient, hence the ladder free space, is isomorphic to a subspace of \(L^1\).  Finally, Lipschitz functions on the two rails extend linearly over the rungs, so the rail free space is complemented in the ladder free space.

\section{The result}

\begin{theorem}
Let
\[
 C=\{(x,0):x\geq 0\}\cup \{(x,x^2):x\geq 0\}
\]
with the metric inherited from the Euclidean plane.  Then \(\F(C)\) is isomorphic to a subspace of an \(L^1\)-space.
\end{theorem}

\begin{proof}
Write
\[
 a(t)=(t,0),\qquad b(t)=(t,t^2)\qquad (t\geq 0).
\]
Let
\[
 P=\{1/n:n\in\N\}\cup \N .
\]
We build a metric graph \(G\).  Its two rails are copies of the two cusp branches.  For each interval
\[
 [1/(n+1),1/n]\quad(n\geq 1),\qquad [m,m+1]\quad(m\geq 1),
\]
put one lower rail edge from \(a(p)\) to \(a(q)\) and one upper rail edge from \(b(p)\) to \(b(q)\), each with its Euclidean chord length.  For each \(p\in P\) add a vertical edge \(V_p\) from \(a(p)\) to \(b(p)\) of length \(p^2\).  Let \(R\subset G\) denote the union of the two rails, with the metric induced from the graph \(G\).

\begin{lemma}
The identity map from the Euclidean cusp \(C\) onto \(R\) is bi-Lipschitz.
\end{lemma}

\begin{proof}
Every edge of \(G\) is realized as a Euclidean segment joining the same endpoints in the plane.  Therefore every graph path between two rail points has length at least the Euclidean distance between those points, so \(d_{\mathrm{Euc}}\leq d_G\) on \(R\).

For points on the lower branch, the graph distance is at most the lower rail length, which is exactly the Euclidean distance.  For points \(b(s),b(t)\) on the upper branch, the graph distance is at most the polygonal length along the parabola.  Since
\[
 \int_s^t (1+4u^2)^{1/2}\,du
 \leq 2\, |s-t|\,(1+(s+t)^2)^{1/2}
 =2\,|b(s)-b(t)|,
\]
the upper-branch graph distance is bounded by twice the Euclidean distance.

It remains to compare cross distances.  Fix \(s,t\geq 0\).  If \(t=0\) there is nothing to prove.  Choose \(p\in P\) with
\[
 |p-t|\leq 2t^2\quad (0<t\leq 1),\qquad |p-t|\leq 1\quad(t\geq 1).
\]
Such a point exists by the definitions of \(P\).  Also \(p\leq 3t\) for \(0<t\leq 1\) and \(p\leq 2t\) for \(t\geq 1\).  The graph path from \(a(s)\) to \(b(t)\) going along the lower rail to \(a(p)\), then across \(V_p\), then along the upper rail to \(b(t)\), has length at most
\[
 |s-p|+p^2+2|p-t|\,(1+(p+t)^2)^{1/2}
 \leq K\bigl(|s-t|+t^2\bigr)
\]
for a universal constant \(K\).  Since
\[
 |a(s)-b(t)|=\bigl((s-t)^2+t^4\bigr)^{1/2}
\]
and \(|s-t|+t^2\leq 2((s-t)^2+t^4)^{1/2}\), this gives \(d_G(a(s),b(t))\leq K' |a(s)-b(t)|\).  The other cross case is symmetric after relabeling the endpoints.  Thus \(d_G\) and the Euclidean cusp metric are bi-Lipschitz equivalent on \(R\).
\end{proof}

Consequently \(\F(C)\) is isomorphic to \(\F(R)\).  We next show that \(\F(R)\) embeds into an \(L^1\)-space.

First, \(\F(R)\) is complemented in \(\F(G)\).  Indeed, if \(f\in\Lip_0(R)\), define \(Ef=f\) on \(R\) and extend \(Ef\) linearly on every vertical edge \(V_p\).  On every rail edge the restriction is already Lipschitz, and on \(V_p\) the endpoint oscillation is at most \(\|f\|_{\Lip}p^2\).  Thus \(E:\Lip_0(R)\to \Lip_0(G)\) is a norm-one pointwise-continuous linear extension operator.  Its preadjoint is a norm-one projection from \(\F(G)\) onto the canonical copy of \(\F(R)\).

It remains to see that \(\F(G)\) is isomorphic to a subspace of \(L^1\).  Orient all edges of \(G\), and let \(L^1(E)\) be the \(L^1\)-sum of the edge intervals with respect to arclength.  For a metric graph, the edge derivative identifies \(\Lip_0(G)\) weak-star continuously with a subspace of \(L^\infty(E)\), and the preadjoint identifies
\[
 \F(G)\simeq L^1(E)/Z,
\]
where \(Z\) is the closed cycle space of the graph.  This is the standard metric-graph form of the Arens--Eells representation.

We show that \(Z\) is complemented in \(L^1(E)\).  Let \(T\) be the spanning tree consisting of:
\begin{itemize}
\item all lower rail edges;
\item all vertical edges \(V_{1/n}\), \(n\geq 1\);
\item all upper rail edges over the intervals \([m,m+1]\), \(m\geq 1\).
\end{itemize}
The chords \(E\setminus T\) are exactly the small upper rail edges over \([1/(n+1),1/n]\) and the large vertical edges \(V_m\), \(m\geq 2\).

For a small upper chord \(e_n=[b(1/(n+1)),b(1/n)]\), its fundamental cycle consists of \(e_n\), the two adjacent vertical rungs, and the lower edge \([a(1/(n+1)),a(1/n)]\).  Since
\[
 1/n-1/(n+1)\asymp 1/n^2
\]
and the upper chord has comparable length, the \(L^1\)-length of this fundamental cycle is bounded by a universal multiple of \(|e_n|\).

For a large vertical chord \(V_m\), \(m\geq 2\), its fundamental cycle consists of \(V_m\), the lower rail path from \(a(1)\) to \(a(m)\), the rung \(V_1\), and the upper rail path from \(b(1)\) to \(b(m)\).  The lower path has length \(m-1\), the upper path is bounded by the parabolic arclength from \(1\) to \(m\), which is \(O(m^2)\), and \(|V_m|=m^2\).  Hence this fundamental cycle also has \(L^1\)-length at most a universal multiple of its chord length.

Let \(c_e\) denote the signed indicator of the fundamental cycle of a chord \(e\), oriented so that \(c_e=1\) on \(e\).  Define
\[
 Qg=\sum_{e\in E\setminus T}\left(\frac1{|e|}\int_e g\,ds\right)c_e,\qquad g\in L^1(E).
\]
The chord supports are disjoint, and the preceding two estimates give
\[
 \|Qg\|_1\leq C\sum_{e\in E\setminus T}\int_e |g|\,ds\leq C\|g\|_1.
\]
Moreover \(Q(c_e)=c_e\) for every fundamental cycle, and the usual fundamental-cycle basis argument gives \(\operatorname{range} Q=Z\).  Thus \(Z\) is complemented in \(L^1(E)\).  Therefore
\[
 L^1(E)/Z \simeq (I-Q)L^1(E),
\]
which is a subspace of \(L^1(E)\).  Hence \(\F(G)\) is isomorphic to a subspace of \(L^1\), and so is its complemented subspace \(\F(R)\).  Since \(\F(C)\simeq\F(R)\), the theorem follows.
\end{proof}

\section{Verification and novelty notes}

The proof is qualitative; no optimized distortion constant is claimed.  The only outside structural input is the standard representation of a metric-graph free space as the edge \(L^1\) space modulo the closed cycle space.  The packet includes the source PDF and a crop of the problem statement.

Bounded novelty check performed before promotion:
\begin{itemize}
\item searched the run registry, solution index, attempt index, and proof-gap index for arXiv `1403.6605', the title, Lipschitz-free product keywords, and cusp/\(L^1\) phrases;
\item searched local parsed sources for later citations of Kaufmann's paper and for the exact cusp phrase;
\item searched the web for exact phrases including ``F(Cusp)'' ``L\^{}1'', ``1403.6605'' ``Cusp'' ``Lipschitz-free'', and ``Products of Lipschitz-free spaces'' ``Cusp''.
\end{itemize}
No direct later answer to the cusp question was found in that bounded search.

\section{References}

\begin{thebibliography}{9}
\bibitem{Godard}
A. Godard, \emph{Tree metrics and their Lipschitz-free spaces}, Proc. Amer. Math. Soc. 138 (2010), 4311--4320.

\bibitem{GodefroyKalton}
G. Godefroy and N. J. Kalton, \emph{Lipschitz-free Banach spaces}, Studia Math. 159 (2003), 121--141.

\bibitem{Kaufmann}
P. L. Kaufmann, \emph{Products of Lipschitz-free spaces and applications}, Studia Math. 226 (2015), 213--227, arXiv:1403.6605.

\bibitem{Weaver}
N. Weaver, \emph{Lipschitz algebras}, World Scientific, 1999.
\end{thebibliography}

\end{document}
